\DocumentMetadata{
  pdfversion=2.0,
  pdfstandard=ua-2,
  lang=en-US,
  tagging=on,
  tagging-setup={math/setup=mathml-SE,tabsorder=structure}
}
\documentclass[11pt,letterpaper]{article}
\usepackage[margin=0.68in,includefoot]{geometry}
\usepackage{newcomputermodern}
\usepackage{amsmath,amscd,mathtools}
\providecommand{\square}{\mdlgwhtsquare}
\usepackage[hidelinks]{hyperref}
\hypersetup{
  pdftitle={Combinatorics PhD Exam, May 2013},
  pdfauthor={University of Florida Department of Mathematics},
  pdfsubject={Graduate mathematics examination},
  pdfdisplaydoctitle=true,
  bookmarksopen=true
}
\setcounter{secnumdepth}{0}
\setlength{\parindent}{0pt}
\setlength{\parskip}{0.28em}
\setlength{\emergencystretch}{3em}
\setlength{\leftmargini}{2.15em}
\setlength{\leftmarginii}{2.65em}
\setlength{\leftmarginiii}{3em}
\renewcommand{\labelenumi}{\textbf{\theenumi.}}
\pagestyle{empty}
\raggedbottom
\allowdisplaybreaks
\newenvironment{examproblems}[1]{%
  \begin{enumerate}%
  \setcounter{enumi}{#1}%
  \setlength{\topsep}{0.35em}%
  \setlength{\partopsep}{0pt}%
  \setlength{\itemsep}{0.48em}%
  \setlength{\parsep}{0.18em}%
}{\end{enumerate}}
\newenvironment{examsubparts}{%
  \begin{enumerate}%
  \setlength{\topsep}{0.18em}%
  \setlength{\partopsep}{0pt}%
  \setlength{\itemsep}{0.16em}%
  \setlength{\parsep}{0.1em}%
}{\end{enumerate}}
\newenvironment{examinstructions}{%
  \par\begingroup\itshape
}{%
  \par\endgroup\vspace{0.18em}
}
\makeatletter
\renewcommand\section{\@startsection{section}{1}{0pt}%
  {-1.25ex plus -.25ex minus -.1ex}%
  {0.55ex plus .1ex}%
  {\normalfont\large\bfseries}}
\renewcommand{\@maketitle}{%
  \newpage\null\vskip -1.2em%
  {\footnotesize\bfseries
    \href{https://www.ufl.edu/}{University of Florida}, \href{https://math.ufl.edu/}{Department of Mathematics}\par}%
  \vskip 0.18em%
  \tagstructbegin{tag=Title}%
    {\LARGE\bfseries\href{https://gma.math.ufl.edu/exams/combinatorics-phd/}{Combinatorics PhD Exam}, May 2013\par}%
  \tagstructend%
  \vskip 0.35em%
  \tagmcbegin{artifact}%
    \hrule height 1pt%
  \tagmcend%
  \vskip 0.55em%
}
\makeatother
\title{Combinatorics PhD Exam, May 2013}
\author{}
\date{}
\begin{document}
\maketitle
\thispagestyle{empty}
\begin{examinstructions}
Answer all questions.
\end{examinstructions}
\begin{examproblems}{0}
\item
State and prove the Erdős–Szekeres theorem on monotone subsequences in permutations.
\item
The 12 edges of a cube are each colored one of two colors, and two such colored cubes are considered equivalent if one can be rotated so that its edge colors agree with the other. Compute the number of equivalence classes of these colored cubes.
\item
How many monic polynomials of degree~\(n\) are there over \(\mathbb{Z}_p[x]\)? How many have no roots in \(\mathbb{Z}_p\)? (Note that two polynomials are considered to be the same if they have the same coefficients, and these coefficients are elements of \(\mathbb{Z}_p\).) Hint: Use inclusion-exclusion.
\item
State and prove Hall’s Marriage Theorem.
\item
Recall that \(R(k,\ell)\) is the minimum integer~\(n\) such that every graph on~\(n\) vertices contains either a clique (complete subgraph) on~\(k\) vertices or an independent set on~\(\ell\) vertices. Prove that \(R(k,\ell)\) is finite for all~\(k\) and~\(\ell\).
\item
Prove that \(R(k,k)>2^{k/2}\) for all~\(k\). Hint: Consider a random graph.
\item
Let \(g_n\) be the number of permutations of length~\(n\) in which each cycle is of length three or longer. Find the exponential generating function of the numbers \(g_n\). Then, determine \(\lim_{n\to\infty} g_n/n!\).
\item
Let \(t_n\) be the number of rooted plane trees on~\(n\) unlabeled vertices in which no vertex has more than two children. Find the ordinary generating function of the numbers \(t_n\). Then determine the exponential growth rate of the sequence \(t_n\). (Just to verify your understanding of the definition of \(t_n\), check that \(t_0=0\), \(t_1=t_2=1\), \(t_3=3\), and \(t_4=9\).)
\end{examproblems}
\end{document}
